\chapter{Introduction}
\label{ch:introduction}

% \chapterepigraph{Architecture begins where engineering ends.}{Walter Gropius}

\openingpara
{Human experience in indoor environments is complex.}
{People do not experience a building only through its physical conditions. They experience it through perception, bodily sensation, emotion, prior experience, social context, ongoing activity, and the wider circumstances of everyday life. A room is shaped not only by its dimensions, materials, temperature, light, or sound, but also by what is taking place within it, who else is present, what its occupants expect, what they remember, and what the space permits or demands of them. An office may feel calm during solitary work and exposed during a difficult conversation. A bright atrium may be experienced as energising, disorienting, or overwhelming. A quiet room may enable concentration while rendering every movement, cough, and fragment of speech conspicuous. Indoor experience therefore cannot only be explained by spatial or environmental conditions alone. It is formed through relations among people, built environments, activities, technologies, and situations.}

At the same time, indoor environments are increasingly shaped by interactive and computational systems. Offices, universities, homes, and public buildings use sensing, automation, feedback, and data-driven control to regulate lighting, temperature, ventilation, access, sound, and other environmental conditions. These systems may be encountered through explicit interfaces, but they may also operate through changes in the environment itself. Buildings are consequently becoming active participants in the conditions through which people work, learn, rest, move, and socialise.\sidenote{I use \emph{interactive indoor environments} to refer to indoor built environments in which spatial, environmental, and computational systems are brought into relation through sensing, automation, adaptation, feedback, or data-driven control.}


\section{the challenge of specifying human experience}

Human--Building Interaction, or HBI, provides the disciplinary context for addressing this development. HBI studies the relations among people, built environments, and interactive technologies. It extends HCI beyond individual devices and interfaces to consider rooms, buildings, environmental systems, infrastructures, and spatial practices as part of interaction. From this perspective, a building is not merely the location in which interaction occurs. Its spatial organisation, environmental conditions, social arrangements, and computational systems contribute to how interaction is perceived and experienced.

Within its research agenda, HBI has articulated two connected concerns. First, human experience in built environments is complex, yet the nature of this complexity remains insufficiently specified. Second, HCI offers theories, methods, and forms of design knowledge that can help investigate this complexity and inform the design of future indoor environments. HBI therefore asks not only how buildings can become interactive, but also what forms of human experience such environments should recognise and support.

The difficulty lies in the breadth of what indoor experience may involve. Existing research has addressed comfort, satisfaction, wellbeing, emotion, bodily experience, atmosphere, aesthetics, social interaction, belonging, participation, inclusion, learning, orientation, and sense of place. Each identifies an important aspect of inhabitation. Their relationship to a broader account of human experience in buildings, however, is not always made explicit. Experience can consequently remain a general term used to acknowledge complexity without specifying what that complexity consists of.

Established approaches to building evaluation provide important but partial accounts. Building science can measure temperature, lighting, sound, air quality, occupancy, and energy use. Comfort and post-occupancy evaluations can assess satisfaction, preference, and perceived environmental quality. These approaches are essential for buildings that are healthy, safe, sustainable, and usable. They are less suited to explaining how environmental conditions acquire meaning within a particular situation.

A space may meet established performance standards while still feeling tiring, exposed, socially uncomfortable, or poorly aligned with the activity taking place. An occupant may describe a room as physically comfortable while also finding it impersonal or difficult to interpret. Another may tolerate environmental discomfort because a location provides privacy, familiarity, or proximity to colleagues. The experiential significance of a condition cannot therefore be inferred from the condition alone.

This produces three connected research problems. The first is conceptual: what forms of human experience should HBI account for? The second is methodological: how can HBI investigate experiences that are contextual, relational, partly ambiguous, and not always easily articulated? The third concerns design: how can richer accounts of experience inform interactive indoor environments without being reduced to additional variables for prediction and optimisation?

The problem addressed in this dissertation is therefore not whether indoor experience is complex. HBI has already recognised that complexity. The problem is how its forms and nuances can be specified, investigated, and made productive for design.

\section{an affective lens for indoor experience}

This dissertation argues that an \emph{affective lens} provides a productive way to address this problem. It demonstrates how affect can help HBI understand the complexity of indoor experience, uncover its forms and nuances, and direct the design of future interactive environments.

Affect is typically approached as an internal state that can be detected, classified, or predicted. Within such approaches, a system may infer whether an occupant is calm, stressed, engaged, or frustrated and adjust its behaviour accordingly. This dissertation takes a different position. It does not treat affect as a stable property of the individual, nor does it propose that buildings should detect and optimise occupants' emotions.

Instead, an affective lens examines how experience is constituted through relations among people, spaces, technologies, activities, and social situations. Emotion is important within this account because it can disclose how a person interprets and relates to their surroundings. Calmness, frustration, exposure, comfort, or attentiveness are not simply responses caused by isolated environmental conditions. They indicate how those conditions have become meaningful within a particular situation.

Calmness, for example, may be associated with relative quiet, daylight, visibility, enclosure, freedom from interruption, or the absence of immediate demands. Its meaning may differ depending on whether an occupant is working, resting, waiting, or observing others. An affective lens retains these relations rather than reducing calmness to a single environmental cause or measurable state.

The thesis demonstrates how this lens can contribute to HBI in two ways. First, it offers a methodological orientation for uncovering the nature and nuances of indoor experiential complexity. Second, it provides a basis for directing and inspiring the design of future interactive environments. The thesis advances HBI through three connected contributions: Affective Interaction as a methodological lens for investigating situated indoor experience; multisensory biophilic design as a design approach, developed through engagement with architecture and spatial expertise, for translating understandings of indoor experience into possibilities for future interactive environments; and an AI-transformed indoor soundscape as a concrete interactive instance through which these methodological and design contributions are brought together and examined.

The argument is developed through a sequence of studies. A scoping review first examines how human experience has been addressed across research on indoor environments. Situated occupant studies then explore the potential of emotion and Affective Interaction for investigating lived experience. The dissertation subsequently shifts from understanding existing environments to imagining future ones through interviews and design sessions with architecture and spatial experts. These studies produce a biophilic design framework for HBI. Finally, an AI-transformed indoor soundscape applies this design direction through a concrete prototype studied using HCI methods.



The overarching research question is:

\begin{quote}
How can an affective lens enable Human--Building Interaction to investigate and specify the complexity of human experience in indoor environments, and inform the design of future interactive indoor environments?
\end{quote}

This question is addressed through four connected research questions:

\begin{itemize}
    \item What forms of human experience are recognised and designed for in research on interactive indoor environments?
    \item What does an affective lens reveal about the complexity of human experience in indoor environments?
    \item How can insights into situated indoor experience inform the design of future interactive indoor environments?
    \item How can one resulting design approach be instantiated and examined through a concrete interactive system?
    \item How can a design approach grounded in human experience be instantiated and examined through a concrete interactive example?
    \item How can the methodological and design insights developed across the thesis be brought together and examined through a concrete interactive system?
\item 
\end{itemize}

\section{mapping human experience design research in interactive indoor environments}

The first part of the dissertation examines how human experience has been addressed in research on interactive indoor environments. A scoping review of 192 papers identifies 11 targeted human experiences, 20 design mechanisms, and seven types of technological intervention. It shows that the field addresses a substantially broader experiential agenda than comfort alone, encompassing emotional, bodily, aesthetic, social, democratic, inclusive, learning, orientational, and place-related experiences.

By specifying this research landscape, the review establishes the conceptual starting point for the dissertation. Emotional experience emerges as a significant strand of research, but its potential as a methodological entry point for investigating situated indoor experience remains less developed. The following study examines this potential by bringing the lens of Affective Interaction from HCI into HBI.

\section{investigating indoor experience through affective interaction}
Building on the potential identified in the scoping review, the next study examines what an affective lens can reveal about lived experience in buildings. It introduces Affective Interaction from HCI as a methodological lens for HBI and applies it through two situated occupant studies: an in-situ survey comparing emotion- and comfort-framed accounts, and a building walk examining how occupants interpret and adjust to their surroundings.

The studies show how occupants sense bodily and environmental conditions, interpret spatial and technological cues, anticipate disruption, and make meaning within social and spatial contexts. From these accounts, the study develops affective practices and corresponding interaction qualities for HBI. It thereby demonstrates how Affective Interaction can extend conventional comfort-oriented methods and provide a richer account of how indoor environments are experienced.


\section{identifying multisensory biophilic design dimensions}
Having examined how an affective lens can reveal the complexity of indoor experience, the dissertation turns from investigation towards design. Because the design indoor experience involves spatial, material, sensory, and technological relations, this stage engages architecture and spatial expertise. Through this engagement, multisensory biophilic design is identified as a promising direction through which these relations can be explored. It concerns how people encounter and relate to natural qualities through space and the senses, while its possibilities within interactive and technologically mediated indoor environments remain underdeveloped.

Through semi-structured interviews with 13 architecture and spatial experts, this study investigates what the legacy, practice, and imaginaries of biophilic architecture can contribute to interactive indoor environments. It identifies eight themes and five design opportunities that extend biophilic interaction beyond visual greenery and restorative effects towards multisensory, interpretive, and ecological experiences. These findings establish a broader design space for multisensory biophilic interaction in future indoor environments.

\section{developing a biophilic interaction design framework}
Building on these expert imaginaries, the following study examines how biophilic possibilities can be articulated through design. Three expert design sessions investigate how architects translate qualities associated with experiences of nature into interactive
spatial propositions for indoor environments. From these sessions, the study develops a biophilic interaction design framework comprising seven biophilic dimensions, seven design values, and six interaction forms. The framework positions biophilic interaction as a multisensory relation among people, technologies, and built environments, and provides the design foundation for the subsequent soundscape study.

\section{an ai-transformed soundscape as a design instance}
The final study brings together the affective and biophilic strands of the dissertation through an AI-transformed indoor soundscape. Informed by the lens of Affective Interaction, the prototype treats indoor experience as situated, relational, and shaped through how occupants sense, interpret, and make meaning within a shared environment. Informed by the biophilic design studies, it explores sound as one sensory and environmentally responsive medium through which biophilic qualities can be introduced indoors.

In contrast to adding static recordings of birds, water, or wind to an office, the prototype transforms ongoing office sound into nature-like soundscapes. An online participant study and expert discussions examine whether this approach can enrich quietness, reframe disruption, and support social, emotional, and biophilic experience, while remaining connected to the activity of the space. The prototype is positioned as a concrete design instance through which the combined affective and biophilic contributions of the dissertation can be examined.


\section{dissertation contributions}
This dissertation advances HBI through three connected contributions.

First, it advances the lens of Affective Interaction as a methodological approach for investigating the complexity of indoor experience. Building on the experiential landscape identified through the scoping review, the occupant studies show how affective inquiry can reveal how people sense, interpret, anticipate, and make meaning within relations among spatial conditions, technologies, activities, and social situations. The resulting Affective Practices and Interaction Qualities provide concepts through which HBI can study and design for lived experience in buildings.

Second, the dissertation develops biophilic interaction as a design approach for technologically mediated indoor environments. Expert interviews broaden the possibilities of biophilic design beyond visual greenery and restorative outcomes, while expert design sessions translate these possibilities into a framework of biophilic dimensions, design values, and interaction forms. Together, these studies provide HBI with a multisensory and spatial design vocabulary for considering relations among people, technologies, built environments, and nature.

Third, the dissertation contributes the design and investigation of an AI-transformed indoor soundscape as a concrete instance of affective and biophilic HBI. The prototype combines an affective understanding of indoor experience with biophilic approaches to multisensory and environmentally responsive design. By transforming ongoing office sound into nature-like auditory textures, the study examines how everyday activity might become material for social, emotional, and biophilic experience within a shared indoor environment.

Together, these contributions position affective HBI as an approach to understanding, imagining, and designing indoor environments from the complexity of human experience. They move from revealing how indoor experience is constituted, to developing design approaches informed by this understanding, and finally to examining their combination through a concrete interactive system.

\section{dissertation structure} 
The remainder of this dissertation is organised as follows.

Chapter~2 presents the scoping review of human-experience design research in buildings, establishing the experiential landscape addressed across smart-building and HBI research.

Chapter~3 introduces the lens of Affective Interaction through situated occupant studies, examining how people sense, interpret, anticipate, and make meaning within indoor environments. 

Chapter~4 draws on interviews with architecture and spatial experts to identify possibilities for multisensory biophilic interaction. 

Chapter~5 develops these possibilities through expert design sessions into a biophilic interaction design framework. 

Chapter~6 brings the affective and biophilic strands together through the design and study of an AI-transformed indoor soundscape.

Chapter~7 synthesises the findings across the dissertation and develops their broader implications along two strands. First, it articulates how the thesis further specifies the future research agenda of HBI, including the forms of human experience, methodological commitments, and design possibilities that require continued investigation. Second, it examines what affective, architectural, and biophilic perspectives on indoor environments contribute to HCI theory, methods, and interaction design. The chapter concludes by identifying unresolved questions and priorities for future research across HBI and HCI.
